% ============================================================================
% CHEATSHEET TEMPLATE
% ============================================================================
% Based on the ANA-G cheatsheet style
%
% Features:
% - 4-column landscape layout on A4
% - Colored boxes for main concepts and sub-concepts
% - Compact formatting optimized for information density
% - Mathematical typesetting support
%
% Usage:
% 1. Replace metadata (title, author, date)
% 2. Adjust number of columns if needed (see \Ncolumn)
% 3. Add your content using mainbox and subbox environments
% 4. Compile with pdflatex (run twice for proper layout)
% ============================================================================

\documentclass[a4paper, fontsize=7pt]{scrartcl}

% ============================================================================
% ENCODING & LANGUAGE
% ============================================================================
\usepackage[utf8]{inputenc}
% \usepackage[german]{babel}    % Uncomment for German hyphenation
% \usepackage[english]{babel}   % Uncomment for English hyphenation

% ============================================================================
% PAGE LAYOUT
% ============================================================================
% 4-column landscape layout with minimal margins
\usepackage[landscape, left=0.75cm, top=1cm, right=0.75cm, bottom=1.5cm, footskip=5pt]{geometry}
\usepackage{flowfram}
\ffvadjustfalse
\setlength{\columnsep}{1cm}

% Number of columns (adjust as needed: 3-4 recommended for A4)
\Ncolumn{4}

% Remove paragraph indentation
\setlength{\parindent}{0cm}

% ============================================================================
% COLORS & BOXES
% ============================================================================
\usepackage[most]{tcolorbox}

% Base style for all boxes
\tcbset{
  base/.style={
    boxrule=0.3mm,
    leftrule=1mm,
    left=1.75mm,
    arc=0mm,
    fonttitle=\bfseries,
    colbacktitle=black!10!white,
    coltitle=black,
    toptitle=0.75mm,
    bottomtitle=0.25mm,
    title={#1}
  }
}

% Main concept box (green accent)
% Use for: definitions, theorems, important formulas
\definecolor{maincolor}{rgb}{0.004, 0.5, 0.35}  % Nice green
\newtcolorbox{mainbox}[1]{
  colframe=maincolor,
  base={#1}
}

% Sub concept box (gray accent)
% Use for: examples, remarks, minor concepts
\newtcolorbox{subbox}[1]{
  colframe=black!20!white,
  base={#1}
}

% Alternative colors (uncomment to use)
% \definecolor{maincolor}{rgb}{0.2, 0.4, 0.8}   % Blue
% \definecolor{maincolor}{rgb}{0.8, 0.2, 0.2}   % Red
% \definecolor{maincolor}{rgb}{0.5, 0.2, 0.6}   % Purple

% ============================================================================
% MATHEMATICS
% ============================================================================
\usepackage{amsmath, amsthm, mathtools, amssymb}
\usepackage{amsfonts, amstext}
\allowdisplaybreaks  % Allow page breaks in equations

% Common math shortcuts (customize as needed)
\def\limn{\lim_{n\to \infty}}
\def\limxo{\lim_{x\to 0}}
\def\limxi{\lim_{x\to\infty}}
\def\limxn{\lim_{x\to-\infty}}
\def\sumk{\sum_{k=1}^\infty}
\def\sumn{\sum_{n=0}^\infty}

% Number sets
\def\N{\mathbb{N}}  % Natural numbers
\def\Z{\mathbb{Z}}  % Integers
\def\Q{\mathbb{Q}}  % Rationals
\def\R{\mathbb{R}}  % Reals
\def\C{\mathbb{C}}  % Complex

% Calculus
\def\dx{\,dx}
\def\dt{\,dt}
\def\dy{\,dy}
\newcommand{\ddx}[1]{\frac{d#1}{dx}}
\newcommand{\ddn}[2]{\frac{d#1}{d#2}}
\newcommand{\pdn}[2]{\frac{\partial #1}{\partial #2}}

% Big O notation
\def\oof{\mathcal{O}}

% ============================================================================
% TABLES
% ============================================================================
\usepackage{tabularx, multirow, booktabs}
\renewcommand*{\arraystretch}{1.5}  % More spacing in tables

% ============================================================================
% LISTS
% ============================================================================
\usepackage{enumitem}
\setitemize{itemsep=0.5pt, topsep=2pt}
\setenumerate{itemsep=0.75pt, topsep=2pt}

% ============================================================================
% GRAPHICS & TIKZ
% ============================================================================
\usepackage{graphicx}
\usepackage{tikz}
\usetikzlibrary{
  positioning,
  calc,
  fit,
  matrix,
  decorations,
  angles,
  quotes,
  babel,
  3d,
  arrows.meta,
  shadings,
  shapes.geometric
}

\tikzset{
  >=stealth,
  % Add custom TikZ styles here
}

% ============================================================================
% HYPERLINKS
% ============================================================================
\usepackage{hyperref}
\hypersetup{
  colorlinks=true,
  linkcolor=maincolor,
  urlcolor=maincolor
}

% ============================================================================
% ADDITIONAL PACKAGES
% ============================================================================
\usepackage{centernot}   % For \centernot command
\usepackage{marvosym}    % Additional symbols
\usepackage{wasysym}     % Additional symbols

% ============================================================================
% DOCUMENT METADATA
% ============================================================================
\title{Cheatsheet ISF}
\author{Thomas Schnyder}
\date{September 2026}

% Table of contents depth
\setcounter{tocdepth}{4}
\setcounter{secnumdepth}{4}

% ============================================================================
% DOCUMENT START
% ============================================================================
\begin{document}

\maketitle

% ============================================================================
% SECTION 1: BEGRIFFE – COMMON GROUND
% ============================================================================
\section{Begriffe -- Common Ground}

\begin{mainbox}{Zeichen -- Daten -- Information -- Wissen}
\begin{itemize}
    \item \textbf{Daten} = Zeichen + Syntax (Form)
    \item \textbf{Information} = Daten + Semantik (Bedeutung)
    \item \textbf{Wissen} = Information + Vernetzung
\end{itemize}
\end{mainbox}

\begin{subbox}{Beispiel Temperatur}
``C2°2'' $\to$ ``22 °C'' $\to$ ``Es ist die aktuelle Aussentemperatur''
$\to$ ``Ich muss eine Jacke anziehen.''
\end{subbox}

\begin{mainbox}{Bedrohungen für Daten \& Informationen}
\begin{itemize}
    \item \textbf{Menschliches Fehlverhalten:} Fahrlässigkeit, Gleichgültigkeit,
          Unwissenheit, Leichtgläubigkeit
    \item \textbf{Organisatorische Schwachstellen:} fehlendes Sicherheitsverständnis
          des Managements, unklare Verantwortlichkeiten, fehlende Prozesse/Strategie/
          Konzepte, mangelhafte Richtlinien \& Awareness, fehlende Kontrollen
    \item \textbf{Technisches Versagen:} ungenügende Wartung, defekte Überwachung (IDS),
          falsch dimensionierte Systeme, Fehlkonfiguration, fehlerhafte
          Software/Firmware/Treiber
    \item \textbf{Höhere Gewalt:} Unwetter, Erdbeben, Feuer, Wasser, Krieg,
          Ausschreitungen
    \item \textbf{Vorsätzliche Angriffe (Internet):} unerlaubter Zugriff, Abhören \&
          Modifizieren, Angriffe auf Verfügbarkeit (DDoS), Systemmissbrauch,
          Viren/Würmer/Trojaner, Drive-by-Infections
\end{itemize}
\textbf{Zufällig} (Naturgewalt, Strom-/TK-Ausfall, Pannen, Bedienerfehler) vs.
\textbf{vorsätzlich} (Malware, Diebstahl, Skript-Kiddies/Hacker,
Wirtschaftsspionage, Missbrauch).
\end{mainbox}

\begin{subbox}{Spektakuläre Cybervorfälle}
\begin{itemize}
    \item \textbf{Stuxnet} (2010): erster Angriff auf SCADA (Supervisory Control And
          Data Acquisition), Siemens-SPS S7, via USB-Stick, Zero-Days.
          $\Rightarrow$ Infrastruktur ist verwundbar, alles ist ``hackbar''.
    \item \textbf{XZ-Utils}: Supply-Chain-Attacke auf OSS (1 Maintainer), Ziel openSSH
          auf Linux; entdeckt von Andres Freund, vermutlich APT29.
          $\Rightarrow$ Lieferketten verwundbar -- weakest link!
    \item \textbf{AI als Hacker} $\Rightarrow$ Automatisierung von Angriffen.
\end{itemize}
Häufiger sind aber alltägliche Fälle: Datendiebstahl, Erpressung, Spionage.
\end{subbox}

\begin{mainbox}{3 Säulen der Informationssicherheit}
\begin{itemize}
    \item \textbf{Technik:} kaufen, konfigurieren
    \item \textbf{Prozesse:} definieren, kontrollieren
    \item \textbf{Mitarbeitende:} sensibilisieren, ausbilden
\end{itemize}
\end{mainbox}

\subsection{Sicherheit \& Risiko}

\begin{mainbox}{Sicherheit}
Zustand, der frei von unvertretbaren Risiken der Beeinträchtigung ist oder als
gefahrenfrei angesehen wird.

\textbf{Absolute Sicherheit ist weder erreichbar noch bezahlbar!}
\end{mainbox}

\begin{mainbox}{Safety vs. Security}
\begin{itemize}
    \item \textbf{Safety} (funktionale Sicherheit, IEC 61508): schützt
          \emph{Menschen} vor Maschinen; Gefahr \emph{zufällig}.
    \item \textbf{Security} (Informationssicherheit, ISO/IEC 27001): schützt
          \emph{Maschinen/Info} vor Menschen; Angriff \emph{willentlich}.
\end{itemize}
\end{mainbox}

\begin{mainbox}{Informations-, IT-, Cyber-Sicherheit}
\begin{itemize}
    \item \textbf{Informationssicherheit:} Schutz der Information als solche,
          unabhängig vom Medium (Datenträger, Papier, Köpfe).
    \item \textbf{IT-Sicherheit:} Schutz von IKT-Systemen gegen vielfältige Gefahren.
    \item \textbf{Cyber-/Internet-Sicherheit:} Schutz von Internet-basierten bzw.
          -verbundenen Systemen (z.B. Browser) gegen Gefahren aus Netzwerken.
\end{itemize}
\textbf{Abgrenzung} (von Solms): Infosec schützt Information \emph{mit und ohne}
IKT-Mittel; Cybersecurity schützt zusätzlich \emph{Nicht-Informationen} mit IKT
(Akteure im Cyberspace und deren darüber erreichbare Assets).
\end{mainbox}

\begin{mainbox}{Informationssicherheit vs. Datenschutz}
\begin{itemize}
    \item \textbf{Informationssicherheit:} Schutz \emph{aller} Daten \& Infos durch
          Massnahmen; Regeln: ISO/IEC 27001. Schutzobjekt: Geschäftsdaten;
          Angreifer will \$\$\$.
    \item \textbf{Datenschutz:} Schutz \emph{personenbezogener} Daten; basiert auf
          informationeller Selbstbestimmung; Regeln: DSG (CH), Branchenvorgaben.
          Schutzobjekt: PII, PHI; Angreifer will \emph{dich}.
\end{itemize}

\textbf{Ziel Infosec:} Bedrohungen \& Schwachstellen identifizieren, bewerten und
Eintrittswahrscheinlichkeit \& Auswirkung von Vorfällen minimieren.
\end{mainbox}

\begin{subbox}{Integrale Sicherheit}
Umfassende Betrachtung \emph{aller} Sicherheitsaspekte einer Organisation --
Informationssicherheit ist nur ein Teil davon.
\end{subbox}

\begin{mainbox}{Risiko}
Möglicher negativer Ausgang eines Vorhabens, verbunden mit Nachteil, Verlust, Schaden.
$$\boxed{\begin{aligned}\text{Risiko} = {} & \text{Eintrittswahrsch.}\\ & \times \text{Schadensausmass}\end{aligned}}$$
Sicherheit und Risiko sind voneinander abhängig: mehr Sicherheit $\Rightarrow$ weniger
Risiko.
\end{mainbox}

\subsection{Schutzziele}

\begin{mainbox}{Grundschutzziele: CIA}
\begin{itemize}
    \item \textbf{Vertraulichkeit} (Confidentiality): Informationen können nicht durch
          unautorisierte Personen, Instanzen oder Prozesse eingesehen werden.
    \item \textbf{Integrität} (Integrity): Daten/Systeme können nicht unautorisiert
          oder zufällig manipuliert bzw. verändert werden.
    \item \textbf{Verfügbarkeit} (Availability): Zugriff auf Dienste/Informationen ist
          in der vom Benutzer gewünschten Zeit möglich.
\end{itemize}
Sicherheitsbetrachtungen müssen sich \emph{mindestens} auf CIA beziehen.
\end{mainbox}

\begin{mainbox}{Wann wird eine Verletzung entdeckt?}
\renewcommand*{\arraystretch}{1.3}
\begin{center}
\begin{tabular}{l|l}
Verfügbarkeit & unmittelbar \\
Integrität, Authentiz. & bald \\
Vertraulichkeit & spät oder nie \\
\end{tabular}
\end{center}
\textbf{``Harvest now, decrypt later''} (Y2Q): heute verschlüsselte Daten
speichern, später (Quantencomputer) entschlüsseln.
\end{mainbox}

\begin{mainbox}{Prüfsumme bis Signatur}
\renewcommand*{\arraystretch}{1.3}
\begin{center}
\begin{tabular}{l|c|c|c|c}
 & \textbf{F} & \textbf{I} & \textbf{A} & \textbf{N} \\
\hline
Prüfsumme & X & & & \\
Hash & X & X & & \\
MAC & X & X & X & \\
Dig. Signatur & X & X & X & X \\
\end{tabular}
\end{center}
F = Fehlererkennung, I = Integrität, A = Authentizität,
N = Nicht-Abstreitbarkeit.
Prüfsumme/Hash nur gegen \emph{zufällige} Fehler; MAC/Signatur auch gegen
\emph{böswillige} Akteure.
\end{mainbox}

\begin{subbox}{Beispiel: ISBN-13 Prüfziffer}
Ziffern 1--12 abwechselnd $\times 1$, $\times 3$ gewichten und summieren; Prüfziffer
= Differenz zum nächsten Zehner.

978-3-3-19-9044-2-?: Summe $= 101 \Rightarrow$ Prüfziffer $= 110 - 101 = 9$.
\end{subbox}

\begin{mainbox}{Privatsphäre vs. Datenschutz}
\textbf{Privatsphäre} = Schutz des persönlichen, nicht öffentlichen Lebensbereichs
(vor Beobachtung, Einmischung, Veröffentlichung). Bsp.: Brief-/Kommunikationsgeheimnis,
Schutz vor Überwachung.

\textbf{Datenschutz} = Kontrolle über die eigenen personenbezogenen Daten
(informationelle Selbstbestimmung).
\end{mainbox}

\begin{mainbox}{Anonymität, Anonymisierung, Pseudonymisierung}
\textbf{Anonymität:} Handlung ist beobachtbar, aber keiner bestimmten Person
zuordenbar (Bsp.: Barzahlung, TOR, I2P).

\textbf{Anonymisierung:} Personenbezug \emph{dauerhaft} entfernt; Zuordnung nicht
oder nur mit unverhältnismässigem Aufwand möglich.

\textbf{Pseudonymisierung:} Personenbezug durch Pseudonym ersetzt; Zuordnung mit
getrennt aufbewahrter Zusatzinfo möglich.

\textbf{Anonymität $>$ Pseudonymität}
\end{mainbox}

\begin{subbox}{Beispiel Blutdruck}
\begin{itemize}
    \item Personenbezogen: ``Reto Meier hat 132/85''
    \item Pseudonymisiert: ``Patient 13267 hat 132/85'' (+ Tabelle
          Reto Meier $\leftrightarrow$ 13267)
    \item Anonymisiert: ``Ein Patient hat 132/85''
\end{itemize}
\end{subbox}

\begin{mainbox}{Weitere Sicherheitsziele}
\begin{itemize}
    \item \textbf{Authentizität:} Echtheit, Überprüfbarkeit, Vertrauenswürdigkeit;
          \emph{Authentifikation} = Prüfung einer behaupteten Eigenschaft.
    \item \textbf{Autorisierung:} Berechtigungszuweisung; nach Identifizierung \&
          Authentisierung Zugriff auf Ressourcen erlauben.
    \item \textbf{Zugriffskontrolle} (accountability): Prüfung der Berechtigung nach
          Authentifikation (Haftung, Verantwortung).
    \item \textbf{Zugreifbarkeit} (accessibility): Fähigkeit zum Zugriff, v.a.
          Barrierefreiheit $\neq$ Benutzerfreundlichkeit.
    \item \textbf{Zuverlässigkeit} (reliability): System funktioniert unter
          bestimmten Bedingungen wie definiert; nahe an Verfügbarkeit \& Qualität.
    \item \textbf{Verbindlichkeit / Nicht-Abstreitbarkeit} (non-repudiation):
          rechtsverbindlicher Zusammenhang zwischen integritätsgeprüftem Inhalt und
          authentisierter Identität (z.B. Onlinebestellung).
    \item \textbf{Vollständigkeit:} Teilaspekt der Integrität.
    \item \textbf{Unbeobachtbarkeit:} Verdeckung eines Kommunikationsvorgangs.
    \item \textbf{Wirksamkeit:} gewünschtes Ergebnis erzielen;
          \textbf{Effektivität:} Genauigkeit \& Vollständigkeit davon.
    \item \textbf{Kopierschutz:} nur autorisierte Nutzer dürfen kopieren
          (z.B. Wasserzeichen).
\end{itemize}
\end{mainbox}

\begin{subbox}{Zielkonflikte \& offene Netze}
Ziele können sich widersprechen bzw. ausschliessen:
\begin{itemize}
    \item Authentifikation vs. Anonymität (Kartenzahlung)
    \item Anonymität vs. Verbindlichkeit (rechtl. Nachweis braucht Identität)
\end{itemize}
\textbf{Internet-Kommunikation} braucht: unveränderte Sender-Identität, Integrität,
Vertraulichkeit, Verbindlichkeit (gesendet/empfangen).
Gewichtung je nach Anwendung (Krypto-Geld, E-Voting, KI \& Privatsphäre).
\end{subbox}

\begin{subbox}{Häufige Fehler \& Tipps}
\begin{itemize}
    \item Hash $\neq$ Authentizität: Angreifer kann Hash einfach neu berechnen.
    \item MAC $\neq$ Nicht-Abstreitbarkeit: gemeinsamer Schlüssel, beide Seiten
          könnten die Nachricht erzeugt haben.
    \item Pseudonymisierte Daten sind \emph{noch} personenbezogen.
    \item Safety $\neq$ Security; Datenschutz $\neq$ Privatsphäre.
\end{itemize}
\end{subbox}

% ============================================================================
% SECTION 2: INTERCEPTING PROXIES
% ============================================================================
\section{Intercepting Proxies}

\subsection{HTTP Grundlagen}

\begin{mainbox}{HTTP Verbindungsablauf}
User gibt URL ein $\to$ Browser (Client) sendet \textbf{HTTP Request} $\to$
Webserver sendet \textbf{HTTP Response} $\to$ Browser zeigt Webseite an.
\end{mainbox}

\begin{mainbox}{HTTP Request (Client $\to$ Server)}
\begin{itemize}
    \item \textbf{Request Line:} Methode, Pfad, HTTP-Version
    \item \textbf{Headers:} Host, Content-Type, Cookie, \ldots
    \item \textbf{Body} (optional): Formulardaten, JSON, \ldots
\end{itemize}
\end{mainbox}

\begin{subbox}{Beispiel Request}
\texttt{POST /login HTTP/1.1}\\
\texttt{Host: www.example.com}\\
\texttt{Content-Type: application/}\\
\texttt{\ \ x-www-form-urlencoded}\\
\texttt{Cookie: session=abc123}\\[2pt]
\texttt{username=admin\&password=1234}
\end{subbox}

\begin{mainbox}{HTTP Response (Server $\to$ Client)}
\begin{itemize}
    \item \textbf{Status Line:} HTTP-Version, Status Code, Reason
    \item \textbf{Headers:} Content-Type, Set-Cookie, \ldots
    \item \textbf{Body:} HTML, JSON, Bild, \ldots
\end{itemize}
\end{mainbox}

\begin{subbox}{Beispiel Response}
\texttt{HTTP/1.1 200 OK}\\
\texttt{Content-Type: text/html}\\
\texttt{Set-Cookie: session=xyz789}\\[2pt]
\texttt{<html><body>Login OK</body></html>}
\end{subbox}

\subsection{Forward, Reverse \& Intercepting Proxy}

\begin{mainbox}{Forward Proxy}
Client (intern) $\to$ Forward Proxy $\to$ Server (extern).\\
\textbf{Schützt den Client, nicht den Server.}
\begin{itemize}
    \item Inhaltsfilter/-kontrolle (z.B. Social Media blockieren)
    \item Überwachung \& Logging (Nachvollziehbarkeit, Detektion)
    \item Caching \& Performance
\end{itemize}
\end{mainbox}

\begin{mainbox}{Reverse Proxy}
Client (extern) $\to$ Reverse Proxy $\to$ Server (intern).\\
\textbf{Schützt den Server, nicht den Client.}
\begin{itemize}
    \item Schutz vor Cyberangriffen (verdächtige Requests blockieren)
    \item Load Balancing (Requests auf Server verteilen)
    \item Caching \& Performance
\end{itemize}
\end{mainbox}

\begin{mainbox}{Intercepting Proxy}
Ein \textbf{Forward Proxy}, der den Datenverkehr \textbf{pausieren, mitlesen und
verändern} kann. Arbeitet nach dem Man-in-the-Middle-Prinzip.
\end{mainbox}

\subsection{Man-in-the-Middle \& TLS}

\begin{mainbox}{MitM-Ablauf mit Intercepting Proxy}
\begin{enumerate}
    \item Browser sendet Request an den Proxy statt direkt an den Server.
    \item Proxy empfängt den Request und zeigt ihn dem Tester an.
    \item Tester kann ihn lesen, verändern oder verwerfen.
    \item Proxy leitet den (ggf. veränderten) Request an den Server weiter.
    \item Response läuft ebenfalls durch den Proxy.
    \item Auch die Response kann inspiziert/verändert werden.
\end{enumerate}
\end{mainbox}

\begin{mainbox}{HTTPS: normal vs. TLS Interception}
\textbf{Normales HTTPS} (TLS Ende-zu-Ende): Proxy sieht nur
\begin{itemize}
    \item \checkmark\ Hostname (via SNI, unverschlüsselt im Handshake)
    \item $\times$ Pfad, Headers, Body
\end{itemize}
\textbf{TLS Interception:} Proxy bricht TLS auf, sieht alles:
\begin{enumerate}
    \item Client $\leftrightarrow$ Proxy: eigene TLS-Verbindung (Proxy-Zertifikat)
    \item Proxy $\leftrightarrow$ Server: eigene TLS-Verbindung (echtes
          Server-Zertifikat)
\end{enumerate}
\end{mainbox}

\begin{mainbox}{CA-Zertifikat des Proxys}
Browser muss dem \textbf{Root-CA-Zertifikat des Proxys} vertrauen:
\begin{enumerate}
    \item ZAP generiert beim ersten Start eine eigene CA + Zertifikat.
    \item CA im Browser als \emph{vertrauenswürdige Zertifizierungsstelle}
          importieren.
    \item ZAP erzeugt dynamisch pro besuchter Domain ein Zertifikat, signiert von
          der eigenen CA.
\end{enumerate}
\end{mainbox}

\subsection{OWASP ZAP}

\begin{mainbox}{OWASP Zed Attack Proxy (ZAP)}
Open-Source Intercepting Proxy (\url{www.zaproxy.org}). Funktionen:
\begin{itemize}
    \item HTTP Requests/Responses abfangen
    \item Manuelles Testen (Break/Intercept)
    \item Automatisierte Scans, Fuzzing, Spidering
\end{itemize}
Alternative: \textbf{Burp Suite Community Edition} (PortSwigger).
\end{mainbox}

\begin{subbox}{Setup}
\begin{enumerate}
    \item ZAP installieren
    \item Browser-Proxy auf \texttt{localhost:8080} setzen
    \item ZAP-CA-Zertifikat exportieren
    \item CA-Zertifikat im Browser importieren
    \item Intercepten
\end{enumerate}
\end{subbox}

\begin{subbox}{Häufige Fehler \& Tipps}
\begin{itemize}
    \item Forward $\leftrightarrow$ Reverse nicht verwechseln: Forward schützt
          den \emph{Client}, Reverse den \emph{Server}.
    \item Intercepting Proxy = Forward Proxy (auf der Client-Seite).
    \item Ohne importierte Proxy-CA: Zertifikatswarnungen im Browser, kein
          HTTPS-Intercept.
    \item Auch bei HTTPS ist der Hostname via SNI sichtbar.
\end{itemize}
\end{subbox}

% ============================================================================
% SECTION 3: ISMS & STANDARDS
% ============================================================================
\section{ISMS \& Standards}

\subsection{Managementsysteme \& ISMS}

\begin{mainbox}{Managementsystem}
Systematisches, geplantes Vorgehen zur \textbf{Zielerreichung} mittels einer
\textbf{Leitlinie/Strategie}; verknüpft vorhandene Elemente durch \textbf{Prozesse}.
Umfasst Organisationsstrukturen, Richtlinien, Planung, Verantwortlichkeiten,
Methoden, Verfahren, Prozesse \& Ressourcen (ISO/IEC 27000, 3.2.5).

Moderne Managementsysteme haben immer einen Rückkopplungs- bzw.
\textbf{ständigen Verbesserungsprozess}: \textbf{PDCA}-Zyklus
(Plan -- Do -- Check -- Act).
\end{mainbox}

\begin{mainbox}{Informationssicherheit (Kernziele)}
Aufrechterhaltung von \textbf{Vertraulichkeit, Integrität, Verfügbarkeit} von
Informationen. Nebenziele: Nicht-Abstreitbarkeit, Privatsphäre, Authentizität,
Anonymität, \ldots\ (ISO/IEC 27000, 3.28).

Meist aus Unternehmenssicht: \emph{alle} Infos, die in einer Organisation
gespeichert, verarbeitet, übertragen werden. ``Datensicherheit'' meint meist
personenbezogene Daten (v.a. Kundendaten).
\end{mainbox}

\begin{mainbox}{ISMS -- Vorgehen}
Unternehmensleitbild, IT-Strategie, Compliance $\to$ \textbf{Risikoanalyse}
(Abhängigkeiten erkennen) $\to$ mit Management abstimmen $\to$
\textbf{IS-Leitlinie} $\to$ IS-Konzepte (Organisation, Notfall, Grundschutz,
Umsetzung, Awareness, Weisungen, Audit) $\to$ IS-Prozess (zyklisch).
\end{mainbox}

\begin{mainbox}{Ohne Management-Support geht nichts!}
Sonst: keine Ressourcen (Zeit, Geld), keine Kompetenzen (Befehls- \&
Umsetzungsgewalt), keine Priorität.\\
\textbf{Das Management trägt die Risiken und entscheidet über die Ressourcen!}
\end{mainbox}

\begin{mainbox}{Begriffe: Leitlinie, Richtlinie, \ldots}
\renewcommand*{\arraystretch}{1.2}
\begin{center}
\begin{tabular}{l|l|l}
\textbf{Englisch} & \textbf{Deutsch} & \\
\hline
Policy & Leitlinie & Überblick \\
Standard, Guideline & Richtlinie & $\downarrow$ \\
Procedures & Verfahren & konkret \\
\end{tabular}
\end{center}
\end{mainbox}

\subsection{ISMS-Typen}

\begin{mainbox}{ISIS12 (12 Schritte)}
\begin{enumerate}
    \item \textbf{Leitlinie erstellen} (Ziele, Sicherheitsniveau, Gesamtverantw.
          der Leitung)
    \item MA sensibilisieren
    \item IS-Team aufbauen (ISB, DSB, IT, Anwender, evtl. QM)
    \item IT-Doku erstellen (Rahmen-, Betriebs-, Notfalldoku)
    \item ITSM einführen (Wartung, Änderung, Störungsbeseitigung)
    \item Kritische Anwendungen identifizieren (CIA, MTA, SLA)
    \item IT-Struktur analysieren (Anw.\ $\to$ Systeme; Schutzbedarf vererben)
    \item Sicherheitsmassnahmen modellieren (Bausteine aus Katalog)
    \item Ist-Soll-Vergleich (evtl.\ Reifegradmessung)
    \item Umsetzung planen (Kosten, Priorisierung)
    \item Umsetzen (Verantw., Termine, Ressourcen)
    \item Revision (weitere Zyklen, evtl.\ Zertifizierung)
\end{enumerate}
MTA = max.\ tolerierbare Ausfallzeit.
\end{mainbox}

\begin{subbox}{Schutzbedarf}
Beschreibt, welcher Schutz für Geschäftsprozesse, die verarbeiteten Infos und die
eingesetzte IT \emph{ausreichend und angemessen} ist (BSI Grundschutz-Kompendium).
\end{subbox}

\begin{mainbox}{Weitere ISMS}
\begin{itemize}
    \item \textbf{VdS 10000:} auf \textbf{KMU} zugeschnitten; Mindestanforderungen
          (Organisation, IS-Leitlinie, IS-Richtlinien, MA, IT-Systeme, Netze,
          Cloud, Zugriffsrechte, Backup, \textbf{Risikoanalyse}, Vorfälle).
    \item \textbf{BSI 200-1:} Basis IT-Grundschutz = \textbf{initialer Verzicht auf
          detaillierte Risikoanalyse}; pauschale Gefährdungen statt Einteilung nach
          Schadenshöhe \& Eintrittswahrsch.
    \item \textbf{ISO/IEC 27001:} ``Nur was dokumentiert ist, ist wirklich
          passiert.'' Führung sichern (Mandat, Ressourcen) $\to$ Umfeld erkennen,
          Ziele setzen $\to$ Situation bewerten $\to$ ändern \& messen $\to$
          Führung unterrichten, von vorne.
\end{itemize}
\end{mainbox}

\begin{mainbox}{ISMS im Vergleich}
\renewcommand*{\arraystretch}{1.2}
\begin{tabularx}{\linewidth}{@{}l|X|X@{}}
& \textbf{Pro} & \textbf{Contra} \\
\hline
ISIS12, VdS 10000 & einfach, sehr konkret & Gebühr, nicht intern.\ anerkannt,
skaliert schlecht \\
\hline
BSI 200-1 & umfassend, skaliert, frei zugänglich, Behörden-Pflicht &
$\sim$1000 S.\ Normtext, umständlich \\
\hline
ISO 27001 & umfassend, skaliert, \textbf{Zertifizierung} & Gebühr, braucht
ISMS-Fachwissen \\
\end{tabularx}
\end{mainbox}

\subsection{Normen \& Normungsorganisationen}

\begin{mainbox}{Normen}
Anwendung \textbf{freiwillig}, Zertifizierung freiwillig. \textbf{Aber:} rechtlich
bindend, wenn \textbf{Vertragsbestandteil}; oft Massstab für ``Best Practice''
(z.B.\ Produkthaftung).

\textbf{Standard vs.\ Framework:} Standards sind \emph{spezifisch} und als Best
Practice akzeptiert; Frameworks sind \emph{allgemein}, allgemein angewandte
Praktiken.
\end{mainbox}

\begin{mainbox}{Organisationen}
\begin{itemize}
    \item \textbf{ISO:} 1947, Genf, 167 Länder (CH Gründungsmitglied); technische
          (MP3), klassifikatorische (Ländercodes), Verfahrensstandards (ISO 9000).
    \item \textbf{IEC:} Elektrotechnik/Elektronik; 1906 London, 1948 Genf; Verein
          nach CH-Recht; Nummern 60000--79999. Gemeinsame Normen: Präfix
          \textbf{ISO/IEC}.
    \item \textbf{EN:} Europ.\ Normen (CEN; CENELEC für Elektrotechnik, ETSI für
          Telekom).
    \item Benennung ``aufsteigend'': \textbf{DIN EN ISO/IEC} = dt.\ Norm auf Basis
          europ.\ Norm auf Basis intern.\ Norm.
\end{itemize}
\end{mainbox}

\begin{subbox}{ISO-Normungsprozess}
00 PWI $\to$ 10 NWIP $\to$ 20 WD (Working Draft) $\to$ 30 CD (Committee Draft)
$\to$ 40 DIS $\to$ 50 FDIS (Final Draft) $\to$ \textbf{60 IS} (veröffentlicht)
$\to$ 90 Review $\to$ 95 Rückzug.
\end{subbox}

\begin{subbox}{Kritik an Normen}
Oft von Einzelinteressen geprägt, schwankende Qualität; viele nicht frei
zugänglich $\Rightarrow$ behindert Kontrolle \& Wettbewerb.
\end{subbox}

\begin{mainbox}{Weitere Standards mit Infosec-Bezug}
\begin{itemize}
    \item \textbf{ITIL:} Organisation \& Prozesse des IT-Service-Managements.
    \item \textbf{COBIT:} GRC-Framework (Governance, Risk, Compliance) von ISACA;
          Bindeglied zw.\ Unternehmenssteuerung und IT-Modellen (ITIL, ISO 27001).
    \item \textbf{PCI DSS:} Sicherheit von Karteninhaberdaten.
    \item \textbf{NIST SP 800-30} (Risikomanagement), \textbf{SP 800-53}
          (Massnahmenkatalog), \textbf{NIST-CSF}.
    \item \textbf{Common Criteria} (ISO/IEC 15408): Prüfung \& Bewertung der
          Sicherheit von \emph{IT-Produkten}.
    \item \textbf{FIPS 140:} Anforderungen an Kryptomodule (HW/SW), Level 1
          (validierter Algorithmus) bis 4 (ungeschützte Umgebung); Level 2:
          Tamper-Detection.
\end{itemize}
\end{mainbox}

\subsection{ISO 27k-Familie}

\begin{mainbox}{ISO/IEC 27001 -- ISMS-Anforderungen}
\begin{itemize}
    \item Einführung, Betrieb, Überwachung, Wartung \& Verbesserung eines
          \emph{dokumentierten} ISMS
    \item aus BS 7799-2:2002, ISO-Norm seit 2005
    \item Prozessansatz \textbf{PDCA}
    \item \textbf{Scope} \& \textbf{Statement of Applicability (SoA)} definieren
          Umfang der Zertifizierung
    \item \textbf{Zertifizierbar!}
\end{itemize}
\end{mainbox}

\begin{mainbox}{ISO/IEC 27002 -- Code of Practice}
\begin{itemize}
    \item \textbf{93 Controls} (Massnahmen/Prüfpunkte) in \textbf{4 Domänen}:
          Organisatorisch (37), Personen (8), Physisch (14), Technologisch (34)
    \item Umsetzungsanleitungen, aber wenig Detailgrad
    \item \textbf{Nicht zertifizierbar} (nur ``sollte''-Formulierungen)
    \item gut für \textbf{Grundschutz} (Mindestanforderungen)
\end{itemize}
\end{mainbox}

\begin{mainbox}{Weitere ISO 27k}
\renewcommand*{\arraystretch}{1.15}
\begin{tabularx}{\linewidth}{@{}l|X@{}}
27000 & Überblick \& Begriffe \\
27004 & Messung (Monitoring, Metriken) \\
27005 & IS-Risikomanagement (Framework, \emph{keine} Methode vorgegeben;
setzt 27001/27002 voraus) \\
27014 & Governance of IS \\
27017/18 & Cloud-Controls / PII in Public Cloud \\
27019 & Energieversorgung (Prozesssteuerung) \\
27031 & Business Continuity \\
27032 & Cybersecurity \\
27033 & Netzwerksicherheit (7 Teile) \\
27034 & Applikationssicherheit \\
27035 & Incident Management \\
27799 & Gesundheitswesen \\
\end{tabularx}
\end{mainbox}

\begin{subbox}{ISO 27035 -- Incident Management}
Erkennung, Reporting, Bewertung \& Behandlung von IS-\emph{Vorfällen} und
-\emph{Schwachstellen}; kontinuierliche Verbesserung. Für grosse \& mittlere Org.
\end{subbox}

\subsection{NIST Cybersecurity Framework}

\begin{mainbox}{NIST-CSF}
NIST: US-Behörde (Handelsministerium). Executive Order 13636 (Obama, 2013):
Framework zum Schutz \textbf{kritischer Infrastrukturen}, risikobasiert,
Best Practices. \textbf{v1.0 2014, v2.0 2024}.

\textbf{Komponenten:}
\begin{itemize}
    \item \textbf{Core:} 6 Functions, 23 Kategorien, 106 Subkategorien (Controls)
    \item \textbf{Tiers:} Wirksamkeit messen, 1 (tief) -- 4 (hoch): Partial, Risk
          Informed, Repeatable, Adaptive
    \item \textbf{Profiles:} Ist-Zustand, Soll-Reifezustand, \textbf{Roadmap};
          ausgerichtet an Geschäftsanforderungen, Risikotoleranz, Ressourcen
\end{itemize}
\end{mainbox}

\begin{mainbox}{6 Functions}
\begin{itemize}
    \item \textbf{Govern} (neu v2.0): Strategie, Policies, Rollen
    \item \textbf{Identify:} Asset-Inventar verstehen
    \item \textbf{Protect:} präventive Controls
    \item \textbf{Detect:} Vorfälle erkennen
    \item \textbf{Respond:} auf Vorfall reagieren
    \item \textbf{Recover:} Resilienz \& Wiederherstellung
\end{itemize}
\end{mainbox}

\begin{mainbox}{IKT-Minimalstandard (CH)}
Vom \textbf{NIST-CSF abgeleitet}, auf CH angepasst (BWL). Empfohlen für
Betreiber \textbf{kritischer Infrastrukturen}, nutzbar für alle. 5 Bereiche
(Identifizieren, Schützen, Detektieren, Reagieren, Wiederherstellen),
\textbf{106 Massnahmen}; Assessment-Tool (Excel).\\
Ausblick: \textbf{CSRM} (Cybersicherheits- \& Resilienzmethode) des BACS.
\end{mainbox}

\subsection{BSI IT-Grundschutz}

\begin{mainbox}{BSI \& Grundschutz-Kompendium}
BSI: dt.\ Bundesbehörde für IT-Sicherheit (1991, Bonn).\\
\textbf{BSI-Standards} = Vorgehensweise; \textbf{Kompendium} = Bausteine,
Gefährdungen, Umsetzungshinweise (8 Domänen, $>$80 Bausteine, 816 S.).
\begin{itemize}
    \item \textbf{200-1:} ISMS-Anforderungen, Zielgruppe Management, kompatibel
          zu ISO 2700x, didaktisch gut
    \item \textbf{200-2:} Grundschutz-Vorgehensweise (Aufbau \& Betrieb ISMS,
          Sicherheitskonzept, Massnahmenauswahl)
    \item \textbf{200-3:} Risikoanalyse (nur bei hohem Schutzbedarf, nicht
          abgedeckten Objekten/Szenarien) $\approx$ ISO 27005
    \item \textbf{200-4:} Business Continuity Mgmt (ersetzt 100-4)
          $\approx$ ISO 27031
\end{itemize}
\end{mainbox}

\begin{mainbox}{Idee \& Absicherungsvarianten}
Abläufe/IT-Komponenten sind überall ähnlich $\Rightarrow$ typische Gefährdungen
wiederverwendbar, anpassbar, erweiterbar.
\begin{itemize}
    \item \textbf{Basis-Absicherung:} Einstieg, schnell alle Prozesse mit
          Basismassnahmen; \emph{kein} Risikomanagement
    \item \textbf{Kern-Absicherung:} Fokus auf ``\textbf{Kronjuwelen}''
    \item \textbf{Standard-Absicherung:} empfohlene Vorgehensweise, umfassend
\end{itemize}
Schutzbedarfskategorien: \textbf{normal, hoch, sehr hoch}.
\end{mainbox}

\begin{mainbox}{Anforderungs-Modalverben}
\textbf{MUSS / DARF NUR:} verpflichtend (Basis-Anf.)\\
\textbf{SOLL:} Abweichung mit Begründung möglich (Standard-Anf.)\\
\textbf{KANN:} fakultativ, Best Practice (erhöhter Schutzbedarf)
\end{mainbox}

\begin{subbox}{Grundschutz vs.\ ISO 27001}
Grundschutz-Anforderungen sind \textbf{konkreter} als ISO 27001 Anhang A.
Zutrittskontrolle 1:1; Backup \textbf{1:n} (1 ISO-Anf.\ $\leftrightarrow$ 12
Anf.\ im Baustein CON.3). \textbf{Nutzen:} passende Bausteine wählen, Massnahmen
umsetzen $\Rightarrow$ wenig Aufwand für Sicherheitskonzept.
\end{subbox}

\subsection{Regulierungen EU \& CH}

\begin{mainbox}{EU-Cybersecurity-Regulierungen}
\begin{itemize}
    \item \textbf{CRA} (Cyber Resilience Act): \emph{Produkte} mit digitalen
          Elementen (HW+SW). Security by Design/Default, Risikobewertung,
          \textbf{SBOM}; Schwachstellenmgmt.\ über Lebenszyklus (5 J.), gratis
          Updates; Vorfälle an \textbf{ENISA} melden; Strafen bis 15 Mio.\ € /
          2.5\,\% Umsatz. Software neu als Produkt (OSS ausgenommen).
    \item \textbf{NIS-2:} sicherer Betrieb \emph{wichtiger Infrastrukturen};
          \textbf{Richtlinie} $\to$ nationale Umsetzung (bis Okt.\ 2024). 18 Sektoren
          (11 kritisch, 7 wichtig); ab $>$50 MA oder $>$10 Mio.\ € Umsatz.
          Meldung: \textbf{24\,h} Frühwarnung, \textbf{72\,h} Detailmeldung.
          Lieferketten, höhere Strafen. ISO 27001 hilft.
    \item \textbf{DORA:} EU-weite \textbf{Verordnung} für \emph{Finanzsektor}
          (CH-Pendant: FINMA-RS 2023/1). IKT-Risikomgmt.\ \& Governance, Info
          Sharing, Resilienztests, Third-Party-Risiko; Verantwortung beim Vorstand.
    \item \textbf{CSA} (Cybersecurity Act): Zertifizierungsrahmen (EUCC auf
          Basis CC), Stärkung ENISA.
\end{itemize}
\end{mainbox}

\begin{mainbox}{CH: Meldepflicht (BACS)}
Seit \textbf{1.\,4.\,2025} für Betreiber kritischer Infrastrukturen: Art,
Auswirkung, Massnahmen des Angriffs innert \textbf{24\,h} melden (später
ergänzbar). Busse bis 100k CHF.
\end{mainbox}

\begin{subbox}{Lieferkette als Angriffsvektor}
Upstream: Tier 3 (Rohstoff) $\to$ Tier 2 $\to$ Tier 1 (Vertragspartner)
$\to$ Vendor $\to$ downstream: Distributor $\to$ Retailer $\to$ Consumer.
\end{subbox}

\begin{subbox}{Häufige Fehler \& Tipps}
\begin{itemize}
    \item ISO \textbf{27001} zertifizierbar, \textbf{27002} nicht (nur ``sollte'').
    \item BSI-Grundschutz verzichtet initial auf detaillierte Risikoanalyse;
          200-3 nur bei Bedarf.
    \item NIS-2 = Richtlinie (national umgesetzt), DORA = Verordnung (direkt).
    \item CRA $\to$ Produkte/Hersteller; NIS-2 $\to$ Betreiber.
    \item IKT-Minimalstandard basiert auf NIST-CSF (nicht ISO).
    \item Policy = Leitlinie (nicht Richtlinie!).
\end{itemize}
\end{subbox}

% ============================================================================
% TIPS FOR AVOIDING OVERFULL BOXES
% ============================================================================
%
% 1. Keep formulas concise - split long equations over multiple lines
% 2. Use shorter variable names in inline math
% 3. Abbreviate table headers
% 4. Use align* environment for multi-line equations
% 5. Test compile frequently to catch overflow early
%
% Example of splitting a long formula:
% Instead of: $$a + b + c + d + e = f + g + h + i + j$$
% Use:
% \begin{align*}
% a + b + c + d + e &= f + g + h \\
%                   &\quad + i + j
% \end{align*}
%
% ============================================================================

\end{document}